\section{端口阻抗、Joule 二次型与离散功率守恒}

\subsection{端口基响应}
令端口峰值复电流相量为 $c\in\C^{n_p}$，source-shape 矩阵为 $S(g)$。在 $e^{+i\omega t}$ 约定下
\begin{equation}
B(g)=-i\omega S(g),
\qquad
A_{\rm em}(g)X(g)=B(g),
\end{equation}
于是任意端口电流下的场为
\begin{equation}
E=X(g)c.
\end{equation}
一次 multi-RHS truth solve 因而覆盖同一几何下的完整端口电流空间。

\subsection{被动端口阻抗的符号}
端口电压方向固定为被动符号约定，使正实部表示端口向电磁系统输入的平均功率。对当前 source orientation，场反应电压定义为
\begin{equation}
v_{\rm field}=-S^T E,
\end{equation}
因此
\begin{equation}
\boxed{
Z_{\rm field}(g)=-S(g)^T X(g).
}
\end{equation}
若线圈采用独立的集中 AC resistance model，则总端口阻抗为
\begin{equation}
\boxed{
Z(a,g)=Z_{\rm field}(g)+R_{\rm wire}(a,g).
}
\end{equation}
其中 $R_{\rm wire}$ 为实对角、非负矩阵。该符号定义必须与独立的离散功率恒等式同时验证。

在 reciprocal material 与一致端口定义下，应有
\begin{equation}
Z_{\rm field}^T=Z_{\rm field}.
\end{equation}
对被动系统还必须有
\begin{equation}
\operatorname{Herm}(Z_{\rm field})\succeq0.
\end{equation}
reciprocity 只给出复对称性，并不能自动保证 passivity。

\subsection{总海水体耗散矩阵}
令导电 Hodge 为 $H_\sigma\succeq0$。定义
\begin{equation}
\boxed{
D_{\rm vol}(g)=X(g)^H H_\sigma(g)X(g)\succeq0.
}
\end{equation}
由于全文采用峰值 phasor，任意端口电流下的周期平均体导电损耗为
\begin{equation}
\boxed{
P_{\rm vol}(g,c)=\frac12 c^H D_{\rm vol}(g)c.
}
\end{equation}
$D_{\rm vol}$ 必须由与 Maxwell 方程相同的 loss operator 构造。当前材料模型只有 conductivity loss；若未来加入 dielectric/magnetic loss，则应把 $D_{\rm vol}$ 推广为包含所有内部被动 EM loss 的统一 $D_{\rm int}$，而不是把新增内部耗散错误归到 outward flux。

\subsection{复 Hermitian Joule 模态张量}
对第 $j$ 个 thermal mode，令 $W_j(g)$ 为把同一导电 Hodge 按 thermal test function $\phi_j$ 加权得到的 Hermitian operator，并保持对 heat projection 的离散一致性。定义
\begin{equation}
\boxed{
H_j(g)=X(g)^H W_j(g)X(g).
}
\end{equation}
因此 $H_j=H_j^H$，但当 thermal mode 变号时一般不半正定。

任意复端口电流下的 cycle-averaged 体 Joule modal source 为
\begin{equation}
\boxed{
q_{{\rm vol},j}(g,c)
=\frac12\operatorname{Re}\!\left(c^H H_j(g)c\right).
}
\end{equation}
不能在构造 $H_j$ 时提前取实部。Hermitian 矩阵虚部的反对称部分会参与有相对相位的多端口局部/modal power。

$D_{\rm vol}$ 不能一般由有限个 $H_j$ 简单求和恢复，因为 thermal basis modes 并不构成空间常数函数的 partition of unity。故 $D_{\rm vol}$ 必须作为独立总功率对象保留，而 $H_j$ 负责 reduced thermal source projection。

\subsection{线圈集中电阻热}
线圈不作为 Maxwell volume conductor 时，其损耗由独立集中 resistance model 表示。若第 $p$ 个端口电阻为 $R_p(a,g)$，对应归一化 cell heat distribution 为 $\ell_p^{(Q)}(g)$，则
\begin{equation}
Q_{\rm wire}
=\sum_p\frac12 R_p(a,g)|c_p|^2\ell_p^{(Q)}(g).
\end{equation}
投影到 thermal basis 后得到 $q_{\rm wire}(a,g,c)$。该模型必须与 full-conductor eddy-current loss 明确区分。

\subsection{开放域 outward-loss matrix 必须独立计算}
对开放电磁域，场端口输入功率除了内部体耗散外，还可能通过外边界/PML 以平均 Poynting flux 离开计算域。由于 Maxwell system 对 $c$ 线性，该外流功率可写成
\begin{equation}
P_{\rm out}(g,c)=\frac12 c^H D_{\rm out}^{\rm phys}(g)c,
\qquad
D_{\rm out}^{\rm phys}=D_{\rm out}^{\rm phys,H}\succeq0.
\end{equation}

truth model 中 $D_{\rm out}^{\rm phys}$ 必须从\textbf{独立的} boundary Poynting-flux bilinear form、absorbing-boundary operator 或 PML absorption operator 构造。禁止在 truth audit 中先定义
\begin{equation}
D_{\rm out}:=\operatorname{Herm}(Z_{\rm field})-D_{\rm vol}
\end{equation}
再用它验证功率恒等式，因为这会使验证退化成恒真式。

真正的 truth-level matrix audit 是
\begin{equation}
\boxed{
\operatorname{Herm}(Z_{\rm field})
\approx D_{\rm vol}+D_{\rm out}^{\rm phys},
}
\end{equation}
误差必须落在 linear solve、boundary/PML quadrature、space discretization 与浮点容差内。该矩阵恒等式一旦成立，对所有 complex current $c$ 自动给出
\begin{equation}
\frac12\operatorname{Re}(c^HZc)
\approx
\frac12c^HD_{\rm vol}c
+\frac12c^HD_{\rm out}^{\rm phys}c
+\frac12c^HR_{\rm wire}c.
\end{equation}

若 Physics Gate 证明目标精度下 outward power 可忽略，才允许使用 $D_{\rm out}^{\rm phys}\approx0$。有限 PEC 截断导致数值上 $D_{\rm out}=0$ 本身不构成开放域物理证据。

\subsection{代理层的联合能量约束}
正式 surrogate 输出集合为
\begin{equation}
\boxed{
\{Z_{\rm field}(g),D_{\rm vol}(g),H_1(g),\ldots,H_r(g)\}.
}
\end{equation}
解码后必须满足
\begin{equation}
Z_{\rm field}^T=Z_{\rm field},
\qquad
D_{\rm vol}=D_{\rm vol}^H\succeq0,
\qquad
H_j=H_j^H.
\end{equation}

surrogate 侧允许定义 implied outward-loss matrix
\begin{equation}
\boxed{
\widehat D_{\rm out}^{\rm imp}
=\operatorname{Herm}(\widehat Z_{\rm field})-\widehat D_{\rm vol},
}
\end{equation}
并要求
\begin{equation}
\widehat D_{\rm out}^{\rm imp}\succeq0.
\end{equation}
这是一条 surrogate feasibility constraint，不是 truth outward power 的定义。validation/audit geometry 上必须把 $\widehat D_{\rm out}^{\rm imp}$ 与独立 truth $D_{\rm out}^{\rm phys}$ 比较。

对当前小端口数，可只对 reciprocal impedance 的耗散块 $R=\operatorname{Herm}Z$ 与 $D_{\rm vol}$ 做联合可行投影，使
\begin{equation}
D_{\rm vol}\succeq0,
\qquad
R-D_{\rm vol}\succeq0,
\end{equation}
并保持 reactive symmetric block 不变。该投影属于固定物理层。

\subsection{固定频率适用范围}
上述 passivity 条件针对当前冻结工作频率的 time-harmonic network。若未来把 frequency 本身加入代理输入，逐频率强制 $\operatorname{Herm}Z(\omega)\succeq0$ 仍不足以保证 causal passive broadband network；此时必须满足 matrix-valued positive-real/causality 条件和跨频率解析结构。当前单频理论不作超出该范围的声明。
